\begin{figure}[H]
\centering
\begin{tikzpicture}[>=Latex,scale=0.98]
    % Tren durante el segundo intervalo
    \draw[black,very thick,rounded corners=2pt,fill=gray!8]
        (-4.8,-0.15) rectangle (4.8,1.2);
    \draw[->,blue!75!black,very thick]
        (1.2,1.68) -- (3.7,1.68)
        node[midway,above] {$\vec v$};

    % Posiciones al comienzo de T_1
    \draw[blue!75!black,very thick,fill=blue!10] (-2.35,0.48) circle (0.13);
    \node[blue!75!black,font=\small,align=center] at (-2.35,-0.78)
        {pelota al iniciar\\el intervalo $T_1$};
    \fill[red!75!black] (4.25,0.8) circle (0.1);
    \node[red!75!black,font=\small,align=center] at (4.1,-0.78)
        {pulso reflejado\\en el frente};

    % Recorridos hasta el encuentro
    \coordinate (E) at (1.15,0.64);
    \draw[->,blue!75!black,very thick]
        (-2.18,0.48) -- (E)
        node[midway,below=5pt] {$wT_1$};
    \draw[->,red!75!black,very thick]
        (4.15,0.8) -- (E)
        node[midway,above=5pt] {$cT_1$};
    \draw[black,thick] (E) circle (0.2);
    \fill[black] (E) circle (0.07);
    \node[font=\scriptsize,align=center,fill=white,inner sep=1pt]
        at (1.15,0.18) {encuentro};

    % Separación inicial
    \draw[gray!65,densely dotted] (-2.35,0.65) -- (-2.35,2.02);
    \draw[gray!65,densely dotted] (4.25,0.9) -- (4.25,2.02);
    \draw[<->,black,thick]
        (-2.35,1.88) -- (4.25,1.88)
        node[midway,fill=white,inner sep=2pt]
        {$D=cT_1+wT_1$};
\end{tikzpicture}
\caption{Segundo intervalo visto por Alice. Después de la reflexión, el pulso avanza hacia la izquierda y la pelota continúa hacia la derecha; la suma de ambos recorridos cubre la separación inicial $D$.}
\label{fig:carrera-alice-t1}
\end{figure}
